% ------------------------------------------------
% 2. DESCRIPCIÓN DEL PROBLEMA Y OBJETIVOS
% ------------------------------------------------

\section{Descripción del problema y objetivos}
\label{sec:problema}

\subsection{El problema del agente viajero}
\label{subsec:tsp}

Se tienen \(n\) ciudades \(V=\{1,\ldots,n\}\) y un costo
\(c_{ij}=c_{ji}\ge0\) de viajar entre cada par de ciudades distintas;
en las instancias de este trabajo, \(c_{ij}\) es una distancia. El TSP
simétrico consiste en encontrar una permutación cíclica de las ciudades,
es decir, una \emph{ruta}, cuyo costo total
\[
f(\sigma)=\sum_{k=1}^{n}c_{\sigma_k\sigma_{k+1}},\qquad
\sigma_{n+1}=\sigma_1,
\]
sea mínimo. Como el costo no depende de la ciudad de partida ni del
sentido del recorrido, el número de rutas distintas es
\((n-1)!/2\). La tabla~\ref{tab:explosion} muestra este número para los
tamaños usados en este trabajo.

\begin{table}[H]
\centering
\small
\begin{tabular}{@{}lrr@{}}
\toprule
Instancia & \(n\) & rutas \((n-1)!/2\) \\
\midrule
aleatoria          &  5 & \(12\) \\
aleatoria          &  8 & \(2\,520\) \\
\texttt{burma14}   & 14 & \(3.1\times10^{9}\) \\
\texttt{ulysses16} & 16 & \(6.5\times10^{11}\) \\
\texttt{ulysses22} & 22 & \(2.6\times10^{19}\) \\
\texttt{berlin52}  & 52 & \(7.8\times10^{65}\) \\
\bottomrule
\end{tabular}
\caption{Número de rutas distintas según el número de ciudades.}
\label{tab:explosion}
\end{table}

Con \(n=5\) es posible enumerar todas las rutas, y por eso esa instancia
se usa para calcular la cadena de Markov de forma exacta. Ya con
\(n=22\), en cambio, evaluar una ruta por nanosegundo tomaría más de
ochocientos años, y con \(n=52\) el número de rutas supera con creces
el número estimado de átomos en la Tierra. La exploración exhaustiva es
inviable, y aquí entran los métodos aleatorizados: el recocido simulado
no recorre el espacio de rutas, sino que lo muestrea mediante una
cadena de Markov diseñada para concentrarse en las rutas de bajo costo.

\subsection{Objetivos}
\label{subsec:objetivos}

\subsubsection*{Objetivo general}

Estudiar el recocido simulado aplicado al problema del agente viajero
como una cadena de Markov, demostrando sus propiedades fundamentales y
contrastándolas con simulaciones reproducibles.

\subsubsection*{Objetivos específicos}

\begin{itemize}
    \item Construir la cadena de Metropolis sobre el espacio de rutas con
    la vecindad 2-opt y demostrar que es irreducible, aperiódica y
    reversible respecto de la distribución de Boltzmann.
    \item Analizar el efecto de la temperatura sobre la distribución
    estacionaria y sobre la velocidad de convergencia, medida por la
    brecha espectral, y relacionar esta última con la profundidad
    crítica del paisaje de costos.
    \item Presentar el recocido simulado como una cadena de Markov no
    homogénea y discutir, con el teorema de Hajek, las diferencias entre
    el enfriamiento logarítmico y el geométrico.
    \item Formular la variante del recocido sobre el modelo QUBO del TSP
    y explicar, a partir de la geometría del paisaje de energía, cómo
    influye el parámetro de penalización \(M\).
    \item Implementar la cadena exacta para \(n=5\) y el recocido
    simulado en un notebook reproducible, y verificar con él los
    resultados teóricos en instancias aleatorias y de TSPLIB.
\end{itemize}
